
\begin{activity}{Modeling Property Scenarios:\\Choosing Physical Property Estimation Methods}

\newcommand{\yourplan}{
    \textbf{Your plan}

    \textbf{2--3 Questions for the Geologist}
    \answerbox[6cm]

    \vspace{0.4em}
    \textbf{Preferred method and why}
    \answerbox[4cm]

    \vspace{0.4em}
    \begin{multicols}{2}
    \noindent\textbf{Pros}
    \begin{itemize}
        \setlength{\itemsep}{0.8em}
        \item \rule{0.88\linewidth}{0.4pt}
        \item \rule{0.88\linewidth}{0.4pt}
    \end{itemize}

    \columnbreak

    \noindent\textbf{Cons}
    \begin{itemize}
        \setlength{\itemsep}{0.8em}
        \item \rule{0.88\linewidth}{0.4pt}
        \item \rule{0.88\linewidth}{0.4pt}
    \end{itemize}
    \end{multicols}
}

\section*{Intructions}
 Below you will be provided with a set of scenarios that you as a geophysicist are required to develop a physical property model for.  The point of this exercise is to develop a model that can be used as a reference based on the available information.  In some cases the best and only option may be to run a geophysical survey and model the property from the results, though generally we can use existing samples and datasets to produce an initial property model.

For the discussions, students should be in groups mixed with geologists and physicists.   
 For each scenario:
 \begin{itemize}
    \item draft 2--3 questions that a practicing geophysicist should ask the geologist for additional information about the geologic setting before devising your scenario,
    \item state your \emph{preferred method} to obtain the required property (e.g., direct measurement, drill-core lab data, composition-based estimate, proxy/correlation, thermodynamic modeling, geophysical inversion/velocity--density regressions), and
    \item justify why you chose it over other options, list pros and cons, and write 2--3 key \emph{questions to ask the geologist}. 
 \end{itemize}
 
\textit{Tip: articulate data availability, scale, $P$--$T$ effects, alteration/weathering, and non-uniqueness when comparing methods.}

% \begin{tabularx}{\linewidth}{@{}lXc@{}}
%     \toprule
%     \textbf{Criterion} & \textbf{Descriptor} & \textbf{Points} \\
%     \midrule
%     Method selection & Chooses a method appropriate to data context, scale, and target property & 6 \\
%     Justification (why over others) & Compares alternatives (measurement vs estimation vs proxies); addresses feasibility & 6 \\
%     Pros & Identifies concrete advantages (accuracy, scalability, cost/time, reproducibility) & 4 \\
%     Cons & Identifies limitations (biases, availability, uncertainty, non-uniqueness) & 4 \\
%     Geo questions & 2--3 targeted, high-value questions to constrain geology & 5 \\
%     Clarity & Answers are concise, organized, and use correct units/terminology & 3 \\
%     \midrule
%     \textbf{Total per scenario} & & \textbf{28} \\
%     \bottomrule
% \end{tabularx}

\newpage
\subsection*{Scenario A: Basin Density Model (Resources/Tectonics)}
\noindent\textbf{Motivation/Goal.} You are supporting a basin analysis that informs \emph{water resource assessment} and subsidence related to extraction. Drill cores are available at multiple wells, and a seismic stratigraphic framework defines a layered structure across the basin. The task is to build a depth-dependent \textbf{density model} suitable for forward gravity modeling and basin compaction studies.  How will you use the available samples and geophysical data to develop the density model?

\yourplan{}

\newpage
\section*{Scenario B: Mountain-Belt Cross-Section (Tectonics)}
\noindent\textbf{Motivation/Goal.} You are preparing a balanced structural cross-section across a metamorphic belt to support a tectonic reconstruction a for a large civil engineering project (construction of a hydroelectic dam). Rock types are mapped at the surface; access to samples are sparse and uneven, though the region has been explored for mineral resources previously. You need \textbf{density} for major units to test isostatic plausibility and gravity fit.  How do you produce a density model?

\yourplan{}

\newpage
\section*{Scenario C: Upper-Mantle Density for Isostasy (Geodynamics)}
\noindent\textbf{Motivation/Goal.} You are part of a team trying to better understand subduction dynamics.  A 3-D regional tomography experiment images $V_s$ and $V_p$ within the mantle down to a depth of 410~km, capturing the mantle wedge, subducted plate and remaining upper mantle.  The geodynamic model requires an upper-mantle \textbf{density} model as an input.  What would be your approach to producing this density model?

\yourplan{}

\newpage
\subsection*{Scenario D: Aeromagnetic Susceptibility Model (Mineral/Exploration)}
\noindent\textbf{Motivation/Goal.} For a greenfields \textbf{mineral exploration} program, you have a regional aeromagnetic survey. Your brief is to produce a \textbf{susceptibility model} that distinguishes key prospective units and to flag where remanent magnetization may invalidate induced-only assumptions.

\yourplan{}

\newpage
\subsection*{Scenario E: Heat-Flow Study in Granitic Terrane (Geothermal/Environmental)}
\noindent\textbf{Motivation/Goal.} An energy agency is scoping geothermal potential and subsurface thermal risk for infrastructure. You must estimate \textbf{thermal conductivity, $k(T)$} and \textbf{heat production} for granitic and metasedimentary units to forecast geotherms and transient responses.  A number of pilot drill holes exist in the basin above, but do not interesct the granite.  The only information about the granite comes from a gravity survey.  How do you estimate the physical properties within the basin and granitic basement?

\yourplan{}

\newpage
\subsection*{Scenario F: Geothermal Heat Flow at the Base of the Antarctic Ice Sheet (Geothermal/Glaciology)}
\noindent\textbf{Motivation/Goal.} You are tasked with developing a geothermal heat flow model into the base of an ice sheet to support assessment of the ice sheet stability due to recent climate change. Several coastal nunataks provide limited outcrops with physical specimens and geochemistry, but their small areal extent and clustered distribution offer only modest constraints on subsurface lithology and properties beneath the ice interior.  There is a relatively low-resolution seismic volume for the region and high-resolution gravity data.  How would you develop the necessary \textbf{thermal conductivity} and \textbf{heat production} models?

\yourplan{}

\end{activity}
